\chapter{Business Domain}

\section{Employee Domain}

An \code{Employee} row is the HR record: name, employee number
(organization-unique), department, employment status
(\code{active}/\code{inactive}/\code{terminated}), optional email/phone,
hire date, and (for a terminated employee) a termination date.

\textbf{Creation} (\code{employees.create\_employee}, admin only)
validates the target department belongs to the caller's organization
and creates the row with \code{employment\_status} as given by the
caller (the create form only offers \code{active}/\code{inactive} ---
see \file{app/forms.py}'s comment on why \code{terminated} is
deliberately excluded from that choice list).

\textbf{Editing} (\code{employees.update\_employee}, admin or a manager
restricted to their own managed departments) accepts a fixed set of
fields (name, department, contact info, employment status, weekly
contract hours) --- reassigning department is a plain field edit here,
not a separate action.

\textbf{Termination} (\code{employees.terminate\_employee}, admin only)
is its own dedicated action, not a generic status edit, because it must
set \code{terminated\_on} in the same step to satisfy the
\code{terminated\_status\_matches\_terminated\_on} CHECK constraint.
Terminating an employee never touches their existing shifts or
attendance history --- see \S4.3 for why a shift already assigned to a
since-terminated employee is left alone.

\textbf{Employee/department relationship} is mandatory (an employee
cannot exist without a department) and validated on every write, not
just at creation.

\textbf{Employee/user relationship} is optional and one-to-one (\S3.4)
--- see \S4.2.

\section{User Account Domain}

A \code{User} is the login identity: email, Argon2 password hash,
\code{role}, and an optional link (\code{employee\_id}) to an
\code{Employee} record.

\begin{center}
\resizebox{0.85\textwidth}{!}{%
\begin{tikzpicture}[node distance=0.5cm]
  \node[flownode] (usr) {User (login, role, password hash)};
  \node[smallnode, below left=0.6cm and 0.3cm of usr] (role) {Role: admin / manager / employee};
  \node[smallnode, below right=0.6cm and 0.3cm of usr] (emp) {Employee (0..1, optional link)};
  \node[smallnode, below=1.2cm of emp] (sched) {Schedule};
  \node[smallnode, below=2.0cm of emp, xshift=-1.6cm] (att) {Attendance};
  \node[smallnode, below=2.0cm of emp, xshift=1.6cm] (leave) {Leave};

  \draw[flowarrow] (usr) -- (role);
  \draw[flowarrow] (usr) -- (emp);
  \draw[flowarrow] (emp) -- (sched);
  \draw[flowarrow] (emp) -- (att);
  \draw[flowarrow] (emp) -- (leave);
\end{tikzpicture}%
}
\end{center}

\textbf{Account creation.} There is no self-service sign-up route
anywhere in the codebase --- every \code{User} row is created directly
(there is no \code{create\_user} service function or route at all; the
only path visible in the code is the \code{flask seed overtime-policy}
CLI command's neighbor pattern, and direct database/CLI provisioning).
This is a genuine gap worth calling out: \VERIFY{} how user accounts
are provisioned in a real deployment is not established by the
application code itself.

\textbf{Authentication, role assignment, and what an account can
access} are covered in full in Chapter~2.

\textbf{Account status.} \code{is\_active} (boolean) and
\code{locked\_until} (temporary, from the lockout policy) are the two
account-status mechanisms that exist. There is no separate
"disabled/pending/invited" status enum --- an account is either active
or not.

\section{Scheduling Domain}

A \code{Shift} represents \emph{planned} work: a department, an
optional assigned employee, a start/end time, a break duration, and a
lifecycle status.

\subsection*{Lifecycle}

\begin{center}
\begin{tikzpicture}[node distance=1.4cm]
  \node[flownode, text width=3.2cm] (draft) {draft};
  \node[flownode, text width=3.2cm, right=of draft] (pub) {published};
  \node[flownode, text width=3.2cm, right=of pub] (cancel) {cancelled};
  \draw[flowarrow] (draft) -- node[above,font=\scriptsize]{publish\_shift (requires an assigned employee)} (pub);
  \draw[flowarrow] (draft.south) .. controls +(0,-1) and +(0,-1) .. node[below,font=\scriptsize]{cancel\_shift} (cancel.south);
  \draw[flowarrow] (pub.south) .. controls +(0,-1.6) and +(0,-1.6) .. node[below,font=\scriptsize]{cancel\_shift} (cancel.south);
\end{tikzpicture}
\end{center}

A shift is created as \code{draft}. Only a draft shift can be edited
directly via \code{update\_shift} (department, times, break, notes) ---
once published, its fields are locked except through the dedicated
\code{assign\_employee} (reassignment is allowed even after publishing,
to cover a call-out), \code{publish\_shift}, and \code{cancel\_shift}
actions. \code{publish\_shift} requires an employee already assigned;
\code{cancel\_shift} never deletes the row (history is preserved for
reporting) and works from either \code{draft} or \code{published}.

\subsection*{Overlap Prevention and Employee Assignment}

Assigning an employee to a shift (\code{create\_shift} with an
\code{employee\_id}, or \code{assign\_employee}) runs three checks:

\begin{enumerate}[nosep]
  \item The employee exists in the caller's organization, is in a
        department the caller may act on, and is currently
        \code{active}.
  \item No overlapping, non-cancelled shift already exists for that
        employee (\longcode{availability.shifts\_overlapping}, backed by
        the database's \longcode{ex\_shifts\_employee\_no\_overlap}).
  \item No approved leave (of a leave type with
        \code{blocks\_scheduling = true}) overlaps the shift's time
        range (\longcode{availability.approved\_leave\_overlapping}) ---
        this is a purely application-level check with no matching
        database constraint, since leave and shifts are different
        tables.
\end{enumerate}

\subsection*{Department Relationships and Coverage}

\code{scheduling.coverage\_summary} answers "do we have enough people
scheduled" with a deliberately minimal pair of counts (published shifts
that day vs.\ active employees in the department) --- not a full
staffing model.

\subsection*{Effective Dates and Validation}

Shifts are not effective-dated in the way pay rates/overtime policies
are; \code{business\_date} is simply the local calendar date of
\code{starts\_at} in the organization's timezone (rule A1), computed
once at creation/update time by the service layer, never re-derived by
the database. A shift is capped at 24 hours
(\code{duration\_max\_24\_hours}), and its break cannot consume the
entire shift (\code{break\_minutes\_less\_than\_duration}).

\section{Attendance Domain}

\code{AttendanceEntry} represents \emph{actual} work: who, when they
started, when (if at all) they stopped, and how the record got there.

\subsection*{Check-In}

\code{attendance.clock\_in(scope, employee\_id=None, at=None)}:

\begin{itemize}[nosep]
  \item Defaults to the caller's own employee record; clocking in
        someone else requires admin/manager (a manager only within
        their own departments).
  \item \code{at} defaults to "now"; only admin/manager may set a
        different (backdated, never future) time --- an employee can
        never backdate their own clock-in. A backdated time is bounded
        to at most 90 days in the past (\code{\_MAX\_BACKDATE}).
  \item Attempts to match the new entry to a single unambiguous
        published shift within a 60-minute grace window on either
        side (\code{\_match\_shift}); zero or more than one candidate
        leaves \code{shift\_id} \code{NULL} rather than guessing.
  \item The database's partial unique index
        (\code{uq\_attendance\_entries\_employee\_id\_open}) is the
        actual guarantee against a duplicate clock-in.
\end{itemize}

\subsection*{Check-Out}

\code{attendance.clock\_out(scope, attendance\_entry\_id, at=None)}
closes an open entry, setting \code{ended\_at} and \code{status =
'closed'}. It explicitly rejects an already-\code{closed} entry, and ---
importantly --- rejects a \code{needs\_review} entry outright: only
\code{correct\_entry} (admin/manager, with a mandatory reason) can
resolve one, "since that defeats the purpose of the review flag."

\subsection*{Needs-Review Behavior}

\code{attendance.flag\_stale\_open\_entries(cutoff\_hours=16)} (run via
the \code{flask attendance flag-stale} CLI command, intended to be
scheduled externally --- there is no in-process scheduler) marks every
\code{open} entry older than the cutoff as \code{needs\_review},
\emph{without ever inventing an end time}. This is the entirety of the
"richer attendance state machine": there is no automatic break
tracking, no multi-step state machine beyond
\code{open}/\code{needs\_review}/\code{closed}. Any richer workflow
(explicit break start/end, multiple open sessions) described elsewhere
as an aspiration is \PLANNED, not implemented --- see \S1.2's product
scope table.

\subsection*{Correction}

\code{attendance.correct\_entry} (admin/manager only, mandatory
\code{edit\_reason}) is the \emph{only} way an entry's times may change
after creation. It can set \code{started\_at}, \code{ended\_at}, and
\code{break\_minutes}; setting \code{ended\_at} always closes the
entry (satisfying the CHECK that ties status to \code{ended\_at}
nullability). It is the only mechanism that can resolve a
\code{needs\_review} entry.

\subsection*{Break Handling}

\PARTIAL{}. There is no \code{BreakSession} table and no break
start/end event anywhere in the schema. \code{break\_minutes} is a
single integer column on both \code{Shift} and \code{AttendanceEntry},
subtracted from the total duration when computing worked/scheduled
hours (\S4.7). There is no support for multiple breaks, break-start/end
timestamps, or a "currently on break" status.

\subsection*{Edge Cases Explicitly Handled}

\begin{itemize}[nosep]
  \item \textbf{Overnight entries} are a single row (not split at
        midnight); \code{business\_date} attributes the whole entry to
        its start date.
  \item \textbf{Missing clock-out} is handled by the needs-review flag,
        never by inventing an end time.
  \item \textbf{Duplicate clock-in} is blocked by the DB's partial
        unique index, independent of any application-level check.
  \item \textbf{A closing time before the start time} is blocked by
        the \code{ended\_after\_started} CHECK (translated to a
        friendly error by \code{\_flush\_or\_raise}).
\end{itemize}

\section{Leave Domain}

\subsection*{State Diagram}

\begin{center}
\begin{tikzpicture}[node distance=1.6cm]
  \node[flownode, text width=2.6cm] (pending) {pending};
  \node[flownode, text width=2.6cm, above right=1.0cm and 2.0cm of pending] (approved) {approved};
  \node[flownode, text width=2.6cm, below right=1.0cm and 2.0cm of pending] (rejected) {rejected};
  \node[flownode, text width=2.6cm, right=4.4cm of pending] (cancelled) {cancelled};

  \draw[flowarrow] (pending) -- node[above,sloped,font=\scriptsize]{approve\_leave} (approved);
  \draw[flowarrow] (pending) -- node[below,sloped,font=\scriptsize]{reject\_leave} (rejected);
  \draw[flowarrow] (pending) -- node[below,font=\scriptsize]{cancel\_leave (owner or admin/manager)} (cancelled);
  \draw[flowarrow] (approved) -- node[above,sloped,font=\scriptsize]{cancel\_leave (admin/manager only)} (cancelled);
\end{tikzpicture}
\end{center}

Only these four states exist (\code{status\_valid} CHECK):
\code{pending}, \code{approved}, \code{rejected}, \code{cancelled}.
Notably, an \emph{approved} request can still be cancelled (by
admin/manager only, "the employee's plans changed") --- the request's
\code{decided\_by\_user\_id}/\code{decided\_at} are deliberately left
untouched by \code{cancel\_leave}, so an approved-then-cancelled request
stays distinguishable from a pending-then-cancelled one using
\code{status} alone (see the extended rationale in
\code{app/models/leave\_request.py}).

\subsection*{Approval and Conflict Detection}

\code{leave.approve\_leave} blocks the approval outright (with a
\code{ValidationError} naming the specific conflicting shift ids and
times) if any \emph{published}, non-cancelled shift overlaps the
requested range for that employee. Self-approval is also blocked: if
the approving admin/manager happens to also be the requesting employee
(compared via \code{scope.employee\_id}), the approval is rejected.
\code{leave.reject\_leave} requires a mandatory \code{decision\_note}
(no such requirement for approval).

\subsection*{Authorization}

An employee may create and cancel (while pending) only their own
requests. Admin/manager may act on requests within their scope
(any department for admin, managed departments for a manager),
including approving/rejecting, and may cancel a request in \emph{any}
status.

\subsection*{Leave Balances}

\PLANNED{}. \code{LeaveType} is a catalog only (name, whether it's
paid, whether it requires approval, whether it blocks scheduling) ---
there is no ledger of days accrued or remaining anywhere in the schema
or service layer.

\section{Overtime Domain}

Overtime is computed by \code{app.services.overtime}, a deliberately
pure, database-free module (no I/O at all) so its tiering arithmetic
can be unit-tested in complete isolation.

\subsection*{Configuration}

An organization's overtime rules are effective-dated data
(\code{OvertimePolicy} + child \code{OvertimeTier} rows), never
hardcoded constants. A policy carries a daily threshold, a weekly
threshold, and a \code{week\_start\_day}; its tiers (each scoped
\code{daily} or \code{weekly}) carry a multiplier and an
\code{[from\_hours, to\_hours)} range \emph{beyond} the relevant
threshold.

\subsection*{Calculation}

Confirmed business rules (from the module's own docstring, not to be
re-derived or second-guessed elsewhere in the codebase):

\begin{itemize}[nosep]
  \item Daily overtime: the first 2 hours beyond the daily threshold
        are paid at 1.5$\times$; daily overtime beyond that (hour 3+ of
        daily OT) is paid at 2.0$\times$.
  \item Weekly overtime: hours beyond the weekly threshold are paid at
        1.5$\times$, but \emph{only} hours not already counted as daily
        overtime for that week --- no worked hour is ever paid under
        two multipliers.
  \item No proration of the weekly threshold for a partial week.
\end{itemize}

For one day, given worked hours $H$, daily threshold $T_d$, and daily
tiers $\{(f_i, t_i, m_i)\}$:

\[
  \text{regular\_hours} = \min(H, T_d), \qquad
  \text{overtime\_hours} = \max(0,\, H - T_d)
\]

\code{overtime\_hours} is then split into per-tier buckets by
\code{apply\_tiers}, each bucket priced separately (\S4.7). For a week,
given each day's already-computed regular hours $r_1,\dots,r_7$ and
weekly threshold $T_w$:

\[
  \text{weekly\_overtime\_eligible} = \max\!\left(0,\; \sum_{i=1}^{7} r_i - T_w\right)
\]

This is exactly what \code{compute\_weekly\_overtime} implements ---
summing \emph{regular} hours (already net of daily OT), never total
worked hours, is what prevents double-counting.

\subsection*{Effective Dates and Policy Resolution}

\longcode{overtime.resolve\_policy(organization\_id, on\_date)} finds the
single policy whose \code{[effective\_from, effective\_to]} range
covers \code{on\_date} (enforced unique at the database level by
\longcode{ex\_overtime\_policies\_organization\_no\_overlap}), returning
\code{None} if the organization has no policy configured for that date
--- callers must handle that case explicitly, never assume a default.

\subsection*{Reports and Authorization}

\code{reports.overtime\_summary} (admin/manager) reports per-employee
overtime \emph{hours} for a department/date range, reusing
\code{labor\_cost.range\_cost\_for\_employee}'s already-correct tiering
so an independent, hand-rolled hours calculation could never
accidentally disagree with the money calculation. It reads only
\code{LineItem.hours}/\code{.category} --- never \code{.rate} or
\code{.cost} --- which is what makes this report safe for a manager to
see under rule A4.

\section{Labor Cost Domain}

Labor cost is the money layer built on three already-verified
components: worked seconds (\code{working\_hours}), overtime tiering
(\code{overtime}), and effective-dated Decimal pay rates
(\code{pay\_rates}). It adds no new business rule of its own --- it
only prices what those three already compute.

\subsection*{What "Labor Cost" Means Here}

For each employee, each day, each priced category (regular time, a
specific daily-OT tier, a specific weekly-OT tier), one \code{LineItem}
is produced:

\[
  \text{cost} = \text{round}_{\text{HALF\_UP}}\bigl(\text{hours} \times \text{rate} \times \text{multiplier},\ 2\text{dp}\bigr)
\]

A department or range total is always the \emph{sum of already-rounded
line items}, never a rounding of one combined raw total --- the two
approaches can disagree by a cent, and the module's own test suite has
a dedicated test proving they do in at least one real case.

\subsection*{Rounding}

\code{Decimal} is used throughout (never \code{float}) specifically
because floating-point binary representation cannot exactly represent
most decimal fractions (e.g.\ \$0.10), which would silently accumulate
rounding error across thousands of line items in a payroll-adjacent
calculation --- exactly the class of bug \file{.claude/rules/time-and-money.md}
exists to prevent. Rounding uses \code{ROUND\_HALF\_UP} to two decimal
places, applied once per line item.

\subsection*{Weekly Overtime Reclassification}

The trickiest part of this module: a naive implementation would price
an hour once as "daily regular" and again as "weekly overtime,"
overpaying. \code{range\_cost\_for\_employee} instead
\emph{reclassifies} weekly-OT-eligible hours out of the regular pool
\emph{before} building the daily "regular" line items, walking backward
from the end of the week (by convention, the hours that pushed the week
over its threshold are the ones worked latest in it) so every worked
hour is priced in exactly one line item. Each reclassified hour is
priced at the rate in force on the specific day it actually came from
--- if an hourly rate changes mid-week, different portions of the same
weekly-OT tier can legitimately be priced at different rates. The
result is \emph{prorated}: a weekly-OT line item is only emitted for
the portion of reclassified hours whose source day actually falls
inside the caller's requested date range, so that
$\text{cost}(\text{Mon..Sun}) = \text{cost}(\text{Mon..Wed}) +
\text{cost}(\text{Thu..Sun})$ for any adjacent split of a week --- this
exact invariant has a dedicated test.

\subsection*{Department Aggregation and What Each Role Can See}

\code{labor\_cost.department\_cost\_summary} (admin or manager, manager
restricted to their own department) returns \textbf{only a total plus a
count of unconfigured employees} --- never a per-employee figure, never
a rate. An employee with genuinely worked hours but no configured pay
rate or overtime policy for the range no longer blanks the whole
department's total (a real bug fixed during development, called "Round
A" in the module's docstring); instead they are excluded and counted,
and the route surfaces that count as an explicit warning rather than a
silently wrong total.

\code{labor\_cost.range\_cost\_for\_employee} (the underlying
per-employee, per-line-item detail, including \code{.rate} and
\code{.cost}) is only ever rendered by the admin-only
\route{/labor-cost/employees/<id>} view --- \emph{no other route in the
codebase calls this function and renders its rate/cost fields}, which
is the actual mechanism (not just a convention) keeping this data away
from managers.

\subsection*{Known, Documented Limitation: Paid Leave Is Not Costed}

\code{LeaveType.is\_paid} is stored on the schema but never read
anywhere in \code{labor\_cost.py} or elsewhere. Approved leave --- paid
or not --- currently contributes exactly \$0 to every labor-cost figure
in the system, because labor cost is derived solely from closed
\code{AttendanceEntry} rows. This is documented explicitly in the
module's own docstring as an open follow-up requiring a real business
decision (does paid leave count toward overtime thresholds? at what
rate is a leave hour priced?), not an oversight.

\section{Reports}

\begin{longtable}{@{}p{2.6cm}p{4.6cm}p{3.2cm}p{3.6cm}@{}}
\toprule
\textbf{Report} & \textbf{Purpose} & \textbf{Filters} & \textbf{Source} \\
\midrule
\endhead
Overtime report & Per-employee total overtime hours for a department/
  range. & Department, start/end date & \code{reports.overtime\_summary}
  (via \route{/reports/overtime}) \\
Hours trend & Per-day total worked hours for a department/range, shown
  as a bar comparison. & Department, start/end date &
  \code{reports.hours\_trend} (via \route{/reports/hours-trend}) \\
Labor cost (department) & Aggregate labor cost total for a department/
  range. & Department, start/end date & \code{labor\_cost.department\_
  cost\_summary} (via \route{/labor-cost}) \\
Labor cost (employee, admin only) & Full per-line-item cost breakdown
  for one employee. & Employee (from URL), start/end date &
  \longcode{labor\_cost.range\_cost\_for\_employee} (via
  \route{/labor-cost/\allowbreak employees/<id>}) \\
\bottomrule
\end{longtable}

All four are department- or employee-scoped and date-range-filtered;
none loads an unbounded table (the routes clamp an excessively wide
\route{?start=}/\route{?end=} range to 92 days --- see \S6.2, security
threat model, for why).

\section{Audit Log}

Every privileged or sensitive action in the system --- login
success/failure, pay rate changes, leave decisions, employee/department
changes, shift publish/cancel, attendance corrections --- is recorded
through exactly one function, \code{audit.record}, which stages (but
never itself commits) an \code{AuditLog} row. Every call site is
required to let its own primary-write commit cover the audit entry too,
so a crash between two separate commits can never leave a state change
persisted with no audit trail.

The log is genuinely \textbf{append-only}: no route or service in the
codebase ever issues an \code{UPDATE} or \code{DELETE} against
\code{audit\_logs}. This is enforced by convention (there is no
database-level trigger preventing it), not by a database permission ---
\VERIFY{} whether the production database role used by the application
has been further restricted at the PostgreSQL grant level to make this
literally impossible rather than merely unused by the current code.

\code{changes} is always a small, deliberately non-sensitive JSON
summary (e.g.\ "a rate was set, for whom, effective when") --- never a
raw dump of the sensitive value itself. A pay-rate change never records
the rate; a leave decision never records the free-text
\code{decision\_note} (which may contain a personal/medical
circumstance).

Only an admin may read the log (\route{/audit-log}, paginated 50 rows
at a time, date-range filtered) or its lighter "Recent Activity"
variant (\route{/recent-activity}, unpaginated but capped, with actor
emails resolved for readability) --- enforced independently inside
\code{audit.list\_entries} and \code{reports.recent\_activity}
themselves, not just by the route decorator.
